%% XAI Survey -- Post-2019 Developments
\documentclass[12pt]{article}

% ---------- Packages ----------
\usepackage{amsmath,amssymb}
\usepackage{graphicx}
\usepackage{booktabs}
\usepackage{multirow}
\usepackage{array}
\usepackage{xcolor}
\usepackage{enumitem}
\usepackage[hidelinks]{hyperref}
\usepackage{url}
\usepackage{xspace}

% ---------- Custom commands ----------
\newcommand{\eg}{\textit{e.g.},\xspace}
\newcommand{\ie}{\textit{i.e.},\xspace}
\definecolor{darkblue}{RGB}{0,51,102}

% =====================================================================
\begin{document}

\title{Advances in Explainable Artificial Intelligence Since 2019:\\
       Methods, Applications, and Emerging Challenges}
\author{Jace King, Claude AI\\{}University of Memphis\\{}jking19@memphis.edu}
\date{}
\maketitle

\begin{abstract}
The field of Explainable Artificial Intelligence (XAI) has undergone rapid
transformation since 2019. Driven by widespread deployment of deep learning in
high-stakes domains, growing regulatory pressure, and the emergence of large-scale
foundation models, the research community has produced new techniques, evaluation
frameworks, and theoretical insights. This survey synthesises key advances across
four themes: (1) post-hoc and ante-hoc explanation methods, including counterfactual
generation, concept-based approaches, and prototype learning; (2) explainability for
attention-based architectures and mechanistic interpretability for large language
models (LLMs); (3) XAI for graph neural networks (GNNs); and (4) evaluation
methodology. We also survey domain applications and close with open problems
including adversarial robustness of explanations, privacy tensions, and the
regulatory landscape of the EU AI Act.
\end{abstract}

\noindent\textbf{Keywords:} Explainable AI, Interpretability, Deep Learning,
Transformers, Large Language Models, Graph Neural Networks, Evaluation,
Counterfactual Explanations, Mechanistic Interpretability, EU AI Act

% =====================================================================
\section{Introduction}
\label{sec:intro}

The opacity of modern machine learning models has long been a source of concern.
Deep neural networks, despite their remarkable predictive power, arrive at
decisions through computations not directly interpretable by humans. This tension
is consequential: a credit-scoring model that cannot explain a rejection may
violate consumer-protection law; a medical diagnostic system whose reasoning
cannot be audited may never gain clinical adoption; an autonomous vehicle that
cannot account for an evasive manoeuvre cannot be certified as safe.

\emph{Explainable Artificial Intelligence} (XAI) encompasses the techniques,
frameworks, and evaluation criteria that seek to make AI behaviour understandable
to human stakeholders. The field matured substantially between 2016 and 2018 with
landmark contributions including LIME~\cite{ribeiro2016lime},
SHAP~\cite{lundberg2017shap}, Integrated Gradients~\cite{sundararajan2017integrated},
Grad-CAM~\cite{selvaraju2017gradcam}, and the theoretical foundations of
Doshi-Velez and Kim~\cite{doshivelez2017rigorous} and
Lipton~\cite{lipton2018mythos}. The 2019 survey by Arrieta et
al.~\cite{arrieta2020xai} marked a consolidation point for that first generation.
Rudin's influential call to abandon black-box models in high-stakes settings in
favour of inherently interpretable models~\cite{rudin2019stop} also appeared in
2019 and framed much subsequent debate.

Since 2019, the landscape has shifted in four important ways. \textbf{First},
Transformers and their descendants (BERT, GPT, ViT) have replaced CNNs and RNNs
as the de facto standard, requiring new explainability strategies. \textbf{Second},
foundation models pre-trained on internet-scale corpora have blurred the boundary
between explanation target and explanation medium: an LLM can itself produce
natural-language explanations. \textbf{Third}, the EU AI Act~\cite{euaiact2021}
mandates transparency requirements for high-risk applications, converting XAI from
a research aspiration to a legal obligation. \textbf{Fourth}, evaluation
methodology has received sustained critical attention, clarifying what ``good''
explanations look like.

This survey covers these developments from 2019 through early 2025, identifying
the most influential conceptual threads rather than attempting encyclopaedic
coverage.

% =====================================================================
\section{Advances in Post-Hoc Explanation Methods}
\label{sec:posthoc}

\subsection{Feature Attribution: Axiomatic Consolidation}

The period 2019--2021 saw significant theoretical consolidation of feature
attribution. Sundararajan and Najmi~\cite{sundararajan2020shapley} showed that
SHAP, Integrated Gradients, and several other methods are instances of a more
general Shapley-value framework parameterised by a choice of \emph{baseline},
clarifying previously confusing discrepancies between methods.

Lundberg et al.~\cite{lundberg2020treeshap} introduced \emph{TreeSHAP}, which
computes exact Shapley values for tree ensembles in polynomial time. Widely used
in production systems, it enabled SHAP to become the dominant explanation
framework for tabular data, and introduced the important distinction between
\emph{path-dependent} and \emph{interventional} SHAP.

\subsection{Counterfactual Explanations}

Counterfactual explanations --- answers to ``what would need to change for the
decision to be different?'' --- emerged as a practically compelling paradigm.
Building on Wachter et al.~\cite{wachter2017counterfactual}, the post-2019
literature focused on counterfactuals that are simultaneously \emph{actionable},
\emph{sparse}, \emph{plausible}, and \emph{diverse}. The DiCE
framework~\cite{mothilal2020dice} addressed diversity directly by formulating
generation as an optimisation problem with an explicit diversity penalty. A
notable thread concerns \emph{causal counterfactuals}, which ground changes in
structural causal models to ensure proposals respect the causal structure of the
world~\cite{verma2020counterfactual}.

\subsection{Concept-Based and Prototype Explanations}

Feature attribution methods identify which inputs mattered but not \emph{why} at
a semantic level. Concept-based methods bridge this gap.
\textbf{ACE}~\cite{ghorbani2019ace} extended TCAV~\cite{kim2018tcav} by
automatically discovering concepts through multi-resolution image segmentation and
activation-space clustering, removing the requirement for manually labelled
concept examples. \textbf{ProtoPNet}~\cite{chen2019protopnet} built
interpretability into the architecture itself: the network learns prototypical
parts used as comparison templates, explaining predictions as ``this patch looks
like prototype $k$.'' \textbf{CRAFT}~\cite{fel2023craft} combined concept
extraction with hierarchical activation factorisation, surfacing concepts at
multiple abstraction levels in a single unified framework.

% =====================================================================
\section{Explainability for Graph Neural Networks}
\label{sec:gnn}

GNNs dominate relational data tasks --- molecular property prediction, social
network analysis, knowledge graph reasoning --- but standard feature-attribution
methods designed for grids or sequences do not transfer because the natural
explanation unit is a subgraph, not a scalar feature.

\textbf{GNNExplainer}~\cite{ying2019gnnexplainer} was the first general post-hoc
method, identifying per-prediction subgraphs and node-feature subsets that
maximise mutual information with the original prediction.
\textbf{PGExplainer}~\cite{luo2020pgexplainer} improved on this by training a
separate network to predict edge masks, enabling generalisation across the test
set rather than solving a per-instance optimisation.
\textbf{SubgraphX}~\cite{yuan2021subgraphx} applied Monte Carlo Tree Search to
select subgraphs maximising a Shapley-value criterion, producing connected
subgraphs that are more semantically coherent and easier for domain experts to
interpret.

A recognised limitation across all three is the \emph{faithfulness evaluation
problem}: because the GNN's decision process is opaque, it is difficult to verify
that an identified subgraph is genuinely causally responsible rather than merely
correlated with the prediction.

% =====================================================================
\section{Explainability in the Era of Transformers}
\label{sec:transformers}

\subsection{The Attention Debate}

The self-attention mechanism produces token-level weights that appear to explain
which tokens influenced an output. Jain and Wallace~\cite{jain2019attention}
challenged this, showing attention weights often disagree with gradient-based
importance scores and that adversarial attention patterns can change distributions
without changing predictions. Wiegreffe and
Pinter~\cite{wiegreffe2019attention} gave a more nuanced response, showing
attention \emph{can} serve as faithful explanation under certain conditions. The
community reached a working consensus: raw attention weights are a weak proxy;
reliable attribution requires propagation through the full computation graph.

Abnar and Zuidema~\cite{abnar2020quantifying} introduced \emph{attention rollout},
recursively multiplying attention matrices across layers. Chefer et
al.~\cite{chefer2021transformer} extended this with gradient signals weighted by
LRP-derived relevance rules, consistently outperforming raw attention and rollout
on faithfulness metrics. Probing work~\cite{tenney2019bert,rogers2020bertology}
further showed that BERT's layers recover the classical NLP pipeline in order ---
lower layers encode part-of-speech, middle layers coreference, higher layers
semantic roles --- suggesting attention heads specialise functionally and that
mechanistic analysis at the head level is tractable.

\subsection{LLMs and Self-Explanation}

The emergence of fluent LLMs created a new mode of explanation: prompting the
model to produce a rationale for its own output. Wei et
al.~\cite{wei2022chainofthought} showed that \emph{chain-of-thought prompting}
both explains the answer and contributes to producing it, improving performance on
complex reasoning tasks. However, LLM-generated explanations can be factually
incorrect or inconsistent with internal computation, motivating a sharp distinction
between \emph{plausible} explanations (coherent to humans) and \emph{faithful}
explanations (accurately describing the model's computation)~\cite{jacovi2020faithfully}.

\subsection{Mechanistic Interpretability}

Mechanistic interpretability seeks to reverse-engineer the algorithms implemented
by neural networks at the level of individual weights and circuits. Elhage et
al.~\cite{elhage2021mathematical} established a mathematical framework for
analysing attention circuits in Transformers. Wang et
al.~\cite{wang2022interpretability} applied this to identify a specific five-layer
circuit in GPT-2 Small responsible for indirect object identification. Conmy et
al.~\cite{conmy2023automating} proposed automated circuit discovery (ACDC) to
improve scalability. Bricken et al.~\cite{bricken2023monosemanticity} applied
dictionary learning to decompose transformer MLP layers into sparse, monosemantic
features, addressing the \emph{superposition hypothesis} that networks represent
more features than neurons by encoding them in overlapping directions.

While these case studies are impressive, circuits identified in small models do not
obviously generalise to billion-parameter deployed models --- a central open
question for the programme.

% =====================================================================
\section{Evaluation of Explanations}
\label{sec:eval}

Evaluation is arguably the most consequential unsolved problem in XAI. Without
rigorous evaluation, it is impossible to compare methods, assess progress, or
trust a deployed explanation system.

Jacovi and Goldberg~\cite{jacovi2020faithfully} formalised \emph{faithfulness} for
NLP, arguing it is the primary desideratum for explanations that audit or debug
models, while \emph{plausibility} (agreement with human intuition) matters for
end-user explanations. These desiderata can conflict. \textbf{ROAR} (RemOve And
Retrain)~\cite{hooker2019roar} operationalises faithfulness by removing features
deemed important, retraining, and measuring the performance drop; it is expensive
but more reliable than perturbation metrics that skip retraining. The
\textbf{ERASER} benchmark~\cite{deyoung2020eraser} provides NLP datasets with
both model-extracted and human-annotated rationales, enabling unified measurement
of faithfulness and plausibility.

Human-centred evaluation adds further complexity. Lage et
al.~\cite{lage2019evaluation} found that simpler, more concise explanations improve
decision speed and that explanation complexity interacts non-trivially with user
expertise. Zhou et al.~\cite{zhou2021evaluating} identified three evaluation types
--- application-grounded, human-grounded, and functionally-grounded --- which
frequently disagree, underscoring the absence of a single authoritative paradigm.

% =====================================================================
\section{Domain Applications}
\label{sec:apps}

\textbf{Healthcare.} XAI faces the most intense deployment pressure in medicine.
Tjoa and Guan~\cite{tjoa2021survey} identified three consumer groups --- clinicians,
patients, and regulators --- each with distinct requirements, and noted that
saliency maps are often inconsistent and vulnerable to perturbations, making them
unreliable for clinical use. Reyes et al.~\cite{reyes2020interpretability} argued
that the key gap in radiology is workflow integration: explanation outputs must
conform to radiological reporting lexicons, not just highlight pixels.

\textbf{NLP.} A distinctive challenge is that tokens are not semantically
independent, so masking one token changes the meaning of surrounding context in
ways that interact non-trivially with model behaviour~\cite{danilevsky2020survey}.

\textbf{Computer Vision.} Concept-based and prototype methods have seen growing
adoption for tasks requiring object-level explanations. ProtoPNet and its
extensions consistently achieve competitive accuracy while being more robustly
interpretable than post-hoc attribution on black-box models~\cite{chen2019protopnet}.

\textbf{Autonomous Systems.} Autonomous vehicles require explanations satisfying
both technical (fault diagnosis, safety verification) and socio-legal (liability,
certification) requirements. Atakishiyev et al.~\cite{atakishiyev2021explainability}
identified a need for \emph{contrastive temporal explanations} that account for
counterfactual decisions not taken and the temporal context over which decisions
unfold.

% =====================================================================
\section{Emerging Challenges and Future Directions}
\label{sec:challenges}

\textbf{Adversarial attacks on explanations.} Slack et al.~\cite{slack2020fooling}
demonstrated that a classifier can behave differently when queried by LIME or SHAP
versus a real user, producing misleading explanations that conceal discriminatory
behaviour. Heo et al.~\cite{heo2019fooling} showed adversarial perturbations can
drastically alter an explanation while leaving the prediction unchanged. Together,
these results imply that explanations cannot be trusted without meta-verification
--- itself an active research problem.

\textbf{Privacy.} Shokri et al.~\cite{shokri2021privacy} showed that model
explanations facilitate membership inference attacks, creating a fundamental
tension in privacy-sensitive domains: the transparency regulators demand may expose
the training data of the individuals regulation is meant to protect. Privacy-
preserving XAI based on differential privacy remains an open problem with
unsatisfactory utility-privacy tradeoffs.

\textbf{Scalability to foundation models.} Most XAI methods were developed on
small models. Gradient-based methods are expensive and unclear under emergent
behaviours; LIME and SHAP require many model evaluations, prohibitive for LLMs.
Mechanistic interpretability offers a principled path, but manual circuit analysis
and combinatorial explosion at scale present serious practical barriers.

\textbf{Regulatory compliance.} The EU AI Act~\cite{euaiact2021}, adopted in 2024,
imposes transparency requirements on high-risk AI systems but does not precisely
define what ``meaningful information about the logic of automated decisions''
requires. This creates substantial legal and technical ambiguity that demands
collaboration between legal scholars, policymakers, and AI researchers.

Looking ahead, the most pressing research directions are: causal XAI grounded in
structural causal models; scalable automated circuit discovery for large models;
shared evaluation benchmarks analogous to ImageNet or GLUE; explanation methods
tailored to generative models; and regulatory operationalisation of XAI
requirements.

% =====================================================================
\section{Conclusion}
\label{sec:conclusion}

The period since 2019 has been transformative for XAI. The field has moved from
a small set of feature-attribution methods applied to CNNs and tabular models to a
rich ecosystem addressing a much wider range of architectures, tasks, and
stakeholder needs. The attention debate, the rise of concept-based explanations,
the emergence of mechanistic interpretability, GNN-specific methods, and
increasing rigour in evaluation are all significant advances.

At the same time, the gap between methodological capability and deployment
practice remains wide. Evaluation is inconsistent, explanations can be
adversarially manipulated, privacy and transparency pull in opposite directions,
and the transition to billion-parameter foundation models has outpaced existing
methods. Rudin's challenge~\cite{rudin2019stop} --- to prefer interpretable models
over explained black boxes --- resonates differently in the era of general-purpose
foundation models, where a single opaque model powers thousands of downstream
applications. Progress will require sustained attention not only to method
development but to evaluation, human factors, causal reasoning, and the
socio-technical contexts in which explanations are produced and consumed.

% =====================================================================
\bibliographystyle{elsarticle-num}
\bibliography{xai_survey}

\end{document}
